\documentclass[11pt,a4paper]{article}
\usepackage[margin=0.85in]{geometry}
\usepackage{graphicx}
\usepackage[sort&compress]{natbib}
%\bibliographystyle{..//..//bib/styles/gcb}
\usepackage{xcolor}

\usepackage[hyphens]{url}
\usepackage{caption}
\usepackage{subcaption}

\setlength\parindent{0pt}
\usepackage{xr}

\usepackage{lineno}
%\usepackage{xr-hyper}
\externaldocument{phenoNAmPhylo}
\externaldocument{phenoNAmPhyloSM}

\newcommand{\lr}[1]{line~\lineref{#1}}

%\lr{refline1}

\begin{document}

\renewcommand{\bibname}{References}%names reference list 

Reviewer comments are italicized and our detailed explanation of changes are in regular text. \\

\textbf{Editor Comments to the Author:}\\
\emph{The manuscript has now been evaluated by two experts in the field. We also thank the authors for their patience during the review process. Both reviewers recognize the novelty and potential value of combining functional trait data with controlled experiments to understand phenological responses to environmental cues. Although their overall recommendations differ substantially, I believe the manuscript has the potential to make an important contribution to our understanding of trait-based predictions of plant phenological shifts, provided several concerns are addressed through a comprehensive revision.}\\

We thank the editor for this positive feedback and are glad to hear they believe our manuscript `has the potential to make an important contribution' to our field.\\ 
        
\emph{ In addition to addressing the reviewers' detailed comments, I particularly encourage the authors to focus on the methodological and interpretive concerns raised by Reviewer 2. The robustness of trait–phenology relationships should be evaluated after accounting for differences between trees and shrubs and considering biologically relevant interactions among environmental cues, particularly forcing × photoperiod and chilling × photoperiod. }

We appreciate the Reviewer's insights and have revised both the methods and results to address their concerns. We have revised our analysis to account for the differences between trees and shrubs from \lr{X} to \lr{Y} and interactions among environmental cues from lr{X} to \lr{Y}.

\emph{Although the factorial design allows temperature and photoperiod to be manipulated independently, clarification is needed on whether actual temperatures were monitored across photoperiod treatments and whether day/night temperature cycles were synchronized with light treatments. These factors could influence thermal exposure and estimated photoperiod sensitivity.}

We have revised the methods from \lr{X} to \lr{Y}  to explain how we monitored the actual temperatures of the growth chambers throughout the experiment and how day/night temperature cycles varied from photoperiod and why from \lr{X} to \lr{Y}.\\

\emph{ Importantly, I fully share Reviewer 2's concern that stronger associations between functional traits and photoperiod sensitivity do not necessarily establish photoperiod as a more important driver of phenological shifts than temperature under natural conditions. The authors should therefore strengthen the mechanistic interpretation, appropriately qualify the conclusions and title, and address the limitations of latitudinal comparisons and potential canopy-position effects on trait measurements.}\\

We see the editor and reviewers point and have revised our discussion, strengthening the mechanistic interpretation from\lr{X} to \lr{Y}. We have also revised the title to better reflect these changes and added a discussion of the latitudinal comparisons and potential canopy-position effects to the discussion from \lr{X} to \lr{Y}.  

\textbf{Reviewer 1 Comments to the Author:}\\
 
\emph{I read and reviewed the article “Traits predict forest phenological responses to photoperiod more than temperature”. The authors of this article investigated the relationship between leaf budburst sensitivity (to temperature and photoperiod) and other plant traits, with the overarching goal of relating phenological sensitivity to the extensive existing plant trait literature. They employed a combination of growth chamber twig phenology experiments (providing info about species-level budburst sensitivity) and trait data that came from combined field collections and trait databases. They found that budburst photoperiod sensitivity was the only phenological trait linked to some of the other plant traits, suggesting that phenological sensitivity to temperature is not a strong component of overall trait strategies. Even the photoperiod sensitivity relationships were not fully related to trait syndromes based on the authors’ predictions, suggesting to me that there exist some complicated and perhaps cryptic relationships between phenological sensitivity and general life-history strategies/selective pressures.}\\

\emph{I thought that the study was well-conceptualized and (for the most part) well-executed, though I had some comments and questions about the field collections that were not described very well in the methods. The statistical analysis was likewise sound. The discussion of results was cohesive and insightful, though I made a couple of suggestions about points that might be interesting to add or explore more if the authors want/have room (though I’ll leave it to them to make those decisions). I don’t have any major criticisms or suggestions for the paper beyond what I address in my line comments below.}

We are happy to hear that the review thought our study `well-conceptualized' and generally 'well-executed' with a `cohesive and insightful' discussion of the results. We see the reviewers point regarding the methods and have added additional details from \lr{X} to \lr{Y} and \lr{X} to \lr{Y}, as discussed in more detail below. We appreciate the interesting points outlined below and have addressed them in this revised version. 

\emph{Line comments:}\\

\emph{L157-160: For the sampled leaf traits, I would expect that leaves of shrubs are relatively similar across an individual since they tend to all be at the very periphery of an individual’s canopy and they all tend to be exposed to the same or similar light conditions. However, we know that leaves of mature trees differ substantially depending on whether they are sun vs. shade leaves and (associated with that difference) what height above the ground they are collected from. Were the tree leaves you quantified traits for collected near the ground or closer to the top of the canopy or both? Were the leaves collected at approximately the same height as the twigs used for the growth chamber? How might your decisions regarding these two questions have impacted your results and the interpretations thereof?}\\

We see the reviewers point and have revised the methods to explain how we consistently sample leaves from the lowest level of the canopy using a pole pruner and subsampled leaves from the tips of those branches to standardize light availability (\R{polePruner} to \R{polePrunerb} and \R{polePruner2} to \R{polePruner2b}). We have also added details on our sampling method fro the growth chamber experiment, in which we used the same method and collected samples from the same area of the canopy (\R{polePruner3} to \R{polePruner3b}). These similarities in methods and the area of the canopy sampled should minimize any variation due to sampling design in our results. 

\emph{Also, related to this point, this got me thinking about how different budburst phenology is at the bottom versus the top of the canopy. I couldn’t find anything in the literature from a quick google scholar search (which surprises me), but I know that Yann Vitasse did something similar with tree seedling phenology monitored in treatments where they were raised at different heights above the forest floor (Vitasse, 2013). Would you expect differences in the relationship between acquisitive/conservative leaf traits and acquisitive/conservative phenology depending on where in the canopy observations and samples were collected?}\\

We agree that microclimate conditions across a canopy can affect budburst phenology, and would predict buds with more exposure to sunlight would differ in their leaf traits due to differences in photosynthetic potential and water loss, and likely earlier budburst in situ (although perhaps not in a controlled environment study such as ours). Due to these differences, we standardized our sampling collection (\R{polePruner} to \R{polePrunerb}, \R{polePruner2}, \R{polePruner3} to \R{polePruner3b}) to ensure we only used buds from the lower canopy at the terminus of branches, thus reducing variation due to microclimates across the canopy.

\emph{L163: Where in the canopy were twigs collected? Near to the ground or near to the top of the canopy?}\\

As mentioned above, we added details on where in the canopy twigs were collected from \R{polePruner3} to \R{polePruner3b}.\\

\emph{L179-180: I was not provided with a copy of the Loughnan \& Wolkovich article that is currently in review, so I was not able to evaluate sample collection and other methods for the western study. The 2018 Flynn \& Wolkovich study answered some of my questions (twigs collected within reach of pole pruners, so likely within ~10 m of ground), but not others (e.g., regarding leaf traits, as described in an earlier comment). Is it possible to add additional info here?}\\

We agree with the reviewer that more detail is needed for this paper to stand alone and have revised our methods to include more details on our sample collection methods  (\lr{methods}, \R{polePruner3},\R{methods3}, \R{methods4}). We employed the same methods to collecting our trait data and conducting the growth chamber experiments in both our eastern and western studies and more explicitly state this point on \lr{methods5} and \lr{polePruner3}. \\

\emph{L236: You described that “traits had a strong relationship with photoperiod”, but wouldn’t it be more accurate to say that “traits had a strong relationship to species-level phenological sensitivity to photoperiod”? Or am I interpreting Fig. 2 incorrectly?}\\

We see the reviewers point and have revised our wording to better reflect that our model estimates trait relationships to species-level phenological responses to photoperiod, such as on \lr{wordChoice} in the results and \lr{wordChoice2} in the discussion.\\

\emph{Fig. 3 caption: Related to my above comment for L236, you describe in the caption that previous studies found “relationships between leaf out timing and wood xylem structures”. But if I’m interpreting your models and figures correctly, it seems like you are instead quantifying the relationship between leaf out sensitivity (not timing) and wood xylem structures (and other traits). Could this not account for the perceived differences in results between your study and the literature?}\\

\emph{L286-287: In line with earlier line comments, perhaps it might be worth mentioning differences across canopy height within an individual in addition to the variation across space that you mention.}\\

This is an interesting point, we have added the following sentence on \lr{X} to discuss how differences across canopy height within an individual could vary in how traits and budburst cue responses may vary.

\emph{L311-319: Sure, soil moisture might be important. But it seems like potential for frost-related leaf mortality might also be worth mentioning as a strong selective force here? Or maybe frost risk is counteracted by capacity for a second, post-frost flush, allowing tree leaves to maintain general trait strategies at the cost of a year of reduced whole-plant carbon budget? Either way, it seems worth mentioning.}\\

This is a great suggestion and have added to our discussion (\lr{X}) discussion of how frost risk and second post-frost flushes relate to general trait strategies in wood plants. 

\emph{Reference cited:
Vitasse, Y. (2013). Ontogenic changes rather than difference in temperature cause understory trees to leaf out earlier. New Phytologist, 198(1), 149–155. https://doi.org/10.1111/nph.12130}\\
 
\textbf{Reviewer 2 Comments to the Author}\\
 
\emph{This manuscript investigates an interesting and timely question: whether interspecific variation in spring phenology can be linked to plant functional strategies through species-specific responses to chilling, forcing, and photoperiod. The combination of field-measured functional traits, controlled-environment budburst experiments, and a joint Bayesian modelling framework is a clear strength. The dataset is substantial, including 47 woody species sampled across eight forest populations in eastern and western North America. However, I have some concerns about the methods and conclusions.}\\

We appreciate that the reviewer found our work `interesting' and that it addresses a `timely question'\\

\emph{1. The authors examine whether functional traits predict species-specific responses to chilling, forcing, and photoperiod. However, trees and shrubs differ substantially in several of the focal traits, such as height and diameter, as well as in phenological timing. Therefore, when trees and shrubs are analysed together, the apparent relationships between traits and environmental sensitivity may partly reflect categorical differences in growth form rather than traits themselves. This issue is further complicated by the strongly unbalanced representation of growth forms between the two transects: the eastern transect contains 13 shrub and 15 tree species, whereas the western transect contains 18 shrub and only four tree species. These analysis should be conducted separately within trees and shrubs.}\\

We see the reviewers point and agree that, although we still found strong differences between trees and shrubs in both trait and phenologies in our analysis, that our community assemblages were unbalanced between our two transects and the effects of this could be explored further. We have added additional analyses to the supplementary material in which we model trees and shrubs individually for each trait.

\emph{2. The analysis of trait variation with latitude is difficult to interpret given that each transect contains only four sites and spans a relatively limited latitudinal range. With limited spatial replication, latitude is likely to be strongly confounded with site identity, local environmental conditions, and species composition. In particular, because species composition changes among sites, the observed latitudinal patterns may reflect species turnover rather than within-species shifts in functional traits. In addition, latitudinal variation in functional traits has already been widely studied in previous studies, which limited its strength.}\\

We can see how our discussion of species assemblages was confusing and requires greater clarification. Species composition did not substantially vary among sites within a transect (with exception of one species in our western transect), as such the latitudinal pattern should reflect within-species shifts in functional traits. \\

\emph{3. The Introduction hypothesized that early, acquisitive species should respond less strongly to chilling, forcing, and photoperiod, whereas later, conservative species should show stronger responses to these cues. However, the physiological or ecological basis for this prediction is not sufficiently developed. In addition, height, diameter, LMA, and leaf nitrogen are associated with photoperiod sensitivity, whereas none of the measured traits clearly explains variation in chilling or forcing sensitivity. However, these results has not been discussed in sufficient biological depth. The Discussion should more directly address why functional traits appear to be linked specifically to photoperiod sensitivity, why comparable relationships are absent for thermal sensitivity, and why some observed trait relationships contradict the expected acquisitive–conservative pattern.}\\

We agree that the justification of our hypotheses could be further developed in the introduction.


\emph{4. The controlled-environment experiment uses a fully crossed design with two levels of chilling, forcing, and photoperiod. However, the phenology model appears to treat these cues primarily as additive effects. This assumption deserves further consideration because interactions among chilling, forcing, and photoperiod are biologically plausible and directly relevant to the manuscript’s interpretation. This is particularly important for the conclusion that photoperiod may constrain warming-induced advances in budburst. Such a mechanism would imply that the effect of forcing depends on photoperiod, i.e., a forcing × photoperiod interaction. Likewise, photoperiod sensitivity may depend on the degree of chilling fulfillment. Testing whether the main trait–photoperiod relationships are robust to the inclusion of biologically relevant cue interactions, particularly forcing × photoperiod and chilling × photoperiod would be useful.}\\

\emph{5. The current framing, particularly the title, gives the impression that photoperiod is more important than temperature in determining budburst timing. However, the results support a more specific conclusion: functional traits are more consistently associated with interspecific variation in photoperiod sensitivity than with variation in chilling or forcing sensitivity. This is not equivalent to showing that photoperiod itself has a stronger effect on budburst than temperature. The manuscript should separate these concepts more carefully throughout the Introduction, Results, Discussion, and title.}\\


\end{document}